% explainer-source-hash: 52c32d41dd9ee9308816b0ea13c8f6fd4352d582538ab1ea52c82ccf33c17cf5
% explainer-source-commit: 4226652b05a73b52389e480fbcab2170bbee92a3
% explainer-generated: 2026-07-16
% Regenerate with /explain-all (see WIRING.md). Do not edit this stamp by hand:
% it is how the toolchain knows whether this document still matches the code.
\documentclass[11pt,a4paper]{article}
\input{preamble}

\title{\vspace{-1.5cm}\textbf{Breakout Game}\\[0.3cm]
\large A Deep Dive into a Game Loop Built Without a Game Engine}
\author{Portfolio explainer: exhaustive walkthrough}
\date{}

\begin{document}
\maketitle
\thispagestyle{fancy}

\begin{keyidea}{What this document is}
A complete, line-aware explanation of \file{breakout-game}: 475 lines of Python
across five modules, using nothing but the standard library. It assumes you have
never written a game, never used Python's \code{turtle} module, and do not know
what a ``game loop'' is. Every term is defined the first time it appears.

It is also honest. Running the code turned up five defects that reading it politely
would have missed, two of which the README states the opposite of. They are in
Section~\ref{sec:honest}, and they are the most valuable part of the document.
\end{keyidea}

\tableofcontents
\newpage

\section{What the project is}

Breakout is an arcade game from 1976. A wall of blocks sits at the top of the screen.
A ball bounces around. The player slides a paddle along the bottom to keep the ball in
play. Every block the ball touches is destroyed and scores a point. Miss the ball and
you lose a life; lose three and the game is over. Destroy every block and you win.

This project rebuilds that from scratch, deliberately without a game engine.

\begin{jargon}{Game engine}
A game engine (Unity, Godot, Pygame) is a library that hands you the hard parts of a
game already solved: a render loop, sprite drawing, collision detection between shapes,
input handling, physics. Building \emph{without} one means you write each of those
yourself. That is the entire point of this project: the exercise is the parts an engine
would have hidden.
\end{jargon}

\begin{jargon}{The \code{turtle} module}
\code{turtle} is a drawing library that ships with Python, designed in the 1960s to
teach children programming. You control a cursor (the ``turtle'') that moves and draws.
It is not a game library: it has no sprites, no collision detection, no frame timing.
It does have a window, shapes that can be moved to coordinates, and keyboard callbacks,
which is just barely enough to build a game on if you supply the rest.
\end{jargon}

\subsection{The whole project at a glance}

\begin{figure}[H]
\centering
\begin{forest} dirtree
[breakout-game
  [main.py {\rmfamily\itshape \ \ 134 lines: screen, state machine, the frame loop}]
  [ball.py {\rmfamily\itshape \ \ 102 lines: velocity, movement, wall/paddle collisions, difficulty}]
  [blocks.py {\rmfamily\itshape \ \ 147 lines: wall generation, layout, block collisions, effects}]
  [paddle.py {\rmfamily\itshape \ \ \ 34 lines: position, key bindings, clamping}]
  [scoreboard.py {\rmfamily\itshape \ \ \ 58 lines: score, lives, high-score persistence}]
  [requirements.txt {\rmfamily\itshape \ \ a comment saying there are no dependencies}]
  [highestScore.txt {\rmfamily\itshape \ \ written at runtime; gitignored}]
]
\end{forest}
\caption{Five modules, one responsibility each. The line counts are measured, not quoted.}
\end{figure}

\begin{keyidea}{The one-sentence architecture}
\file{main.py} owns the loop and the state; the other four modules are passive objects
that only act when the loop calls a method on them. Nothing except \file{main.py} knows
that a frame exists.
\end{keyidea}

\subsection{How the modules relate}

\begin{figure}[H]
\centering
\resizebox{\textwidth}{!}{
\begin{tikzpicture}
  \node[box, minimum width=48mm] (main) {\textbf{main.py}\\Screen, state machine, the frame loop};
  \node[extbox, above=12mm of main] (turtle) {Python stdlib: \code{turtle} (on \code{tkinter})};

  % four modules on one row, generously spaced so no label can land on a node
  \node[box, minimum width=30mm] at ([xshift=-65mm]$(main.south)+(0,-30mm)$) (paddle)
    {\textbf{paddle.py}\\\code{Paddle(Turtle)}\\keys, clamping};
  \node[box, minimum width=30mm] at ([xshift=-21mm]$(main.south)+(0,-30mm)$) (ball)
    {\textbf{ball.py}\\\code{Ball(Turtle)}\\dx, dy, collisions};
  \node[box, minimum width=30mm] at ([xshift=21mm]$(main.south)+(0,-30mm)$) (blocks)
    {\textbf{blocks.py}\\\code{Block}\\wall, effects};
  \node[box, minimum width=34mm] at ([xshift=65mm]$(main.south)+(0,-30mm)$) (score)
    {\textbf{scoreboard.py}\\\code{Scoreboard(Turtle)}\\score, lives};
  \node[store, below=14mm of score] (file) {highestScore\\.txt};

  \draw[dflow] (turtle) -- node[lbl,right]{window, shapes, keys} (main);
  \draw[flow] (main.south) -- node[lbl,pos=0.45,left]{\code{.move()}, \code{.check\_*()}} (ball.north);
  \draw[flow] (main.south) -- (paddle.north);
  \draw[flow] (main.south) -- node[lbl,pos=0.45,right]{\code{.check\_collision\_ball()}} (blocks.north);
  \draw[flow] (main.south) -- (score.north);
  \draw[flow] (blocks.south) to[out=-90,in=-90] node[lbl,below]{\code{.increase\_score()}} (score.south);
  \draw[flow] (score.east) to[out=0,in=0] node[lbl,right]{read at import,\\write at game end} (file.east);
  \draw[flow] (ball.south) to[out=-90,in=-90] node[lbl,below]{reads \code{paddle.xcor()}} (paddle.south);
\end{tikzpicture}}
\caption{Component map. Solid arrows are calls; the dashed arrow is a dependency.
Note that \code{blocks.py} calls into \code{scoreboard.py} directly, and \code{ball.py}
reads the paddle's position: those are the only two couplings that do not go through
\file{main.py}.}
\end{figure}

\section{The game loop, the central idea}

\begin{jargon}{Frame}
A single still image of the game. Video is an illusion built from many stills shown in
quick succession; a game draws one image, moves everything a tiny bit, draws the next,
forever. Each cycle is a frame.
\end{jargon}

\begin{jargon}{Game loop}
The \code{while True} loop that never ends until the player quits. Every pass through it
is one frame, and each pass does the same three things: move things, check what hit what,
draw the result.
\end{jargon}

\subsection{Why \code{tracer(0)} is the load-bearing line}

By default \code{turtle} redraws the window immediately after every single command, and
animates the turtle sliding to its destination. For a drawing toy that is friendly. For a
game it is fatal: the ball would visibly crawl, and the screen would tear and flicker.

\file{main.py} line 15 disables that:

\begin{lstlisting}[language=Python]
wn.tracer(0)   # stop redrawing automatically after every command
\end{lstlisting}

With the tracer off, nothing appears on screen at all until the program explicitly asks
for a redraw with \code{wn.update()}, which the loop does exactly once per frame (line 93).

\begin{keyidea}{Why this turns a toy into a game}
\code{tracer(0)} plus one \code{wn.update()} per iteration is the whole trick. It converts
turtle from ``draw each change instantly'' into double-buffered, frame-based rendering:
move everything invisibly, then present one complete image. Every object in the frame
appears to move at the same instant, because it does.
\end{keyidea}

\subsection{The frame loop, exactly as written}

This is the heart of the project. \file{main.py} lines 91 to 133:

\begin{figure}[H]
\centering
\resizebox{0.93\textwidth}{!}{
\begin{tikzpicture}[node distance=7mm,
  flowstep/.style={rectangle, rounded corners=2pt, draw=inkblue, fill=softgrey,
                   text=inkblue, align=center, inner sep=4pt, minimum height=7mm,
                   minimum width=44mm, font=\small}]

  \node[flowstep, fill=white, draw=linegrey] (top) {\code{while True:}};
  \node[flowstep, below=of top] (update) {\code{wn.update()}\\\scriptsize present the frame};
  \node[dec, below=of update, minimum width=20mm] (st) {state ==\\``playing''?};

  \node[flowstep, below=11mm of st] (c1) {\textbf{1.} \code{block.check\_collision\_ball()}\\\scriptsize the first of two};
  \node[flowstep, below=of c1] (mv) {\textbf{2.} \code{ball.move()}\\\scriptsize x += dx, y += dy};
  \node[flowstep, below=of mv] (cw) {\textbf{3.} \code{ball.check\_collision\_wall()}};
  \node[flowstep, below=of cw] (cp) {\textbf{4.} \code{ball.check\_collision\_paddle()}\\\scriptsize also drives the difficulty ramp};
  \node[flowstep, below=of cp] (pw) {\textbf{5.} \code{paddle.check\_collision\_wall()}};
  \node[flowstep, below=of pw, draw=warnred] (c2) {\textbf{6.} \code{block.check\_collision\_ball()}\\\scriptsize the second, identical call};
  \node[flowstep, below=of c2] (fx) {\textbf{7.} \code{block.update\_effects()}};

  \node[dec, below=10mm of fx, minimum width=20mm] (win) {no blocks\\left?};
  \node[dec, below=13mm of win, minimum width=20mm] (over) {\code{ball.check}\\\code{\_game\_over()}?};
  \node[dec, below=15mm of over, minimum width=18mm] (last) {last\\life?};

  \node[flowstep, right=20mm of win, draw=okgreen] (won) {\code{save\_high\_score()}\\\code{show\_end\_screen(won=True)}\\state = ``won''};
  \node[flowstep, below=11mm of last, draw=warnred] (go) {\code{save\_high\_score()}\\\code{show\_end\_screen(won=False)}\\state = ``game\_over''};
  \node[flowstep, right=22mm of last] (miss) {\code{ball.reset\_after\_miss()}\\\scriptsize keep speed, re-centre};

  \node[flowstep, left=30mm of st, draw=linegrey] (idle) {\code{block.update\_effects()}\\\scriptsize only; no physics};

  \draw[flow] (top) -- (update);
  \draw[flow] (update) -- (st);
  \draw[flow] (st) -- node[lbl,right]{yes} (c1);
  \draw[flow] (st) -- node[lbl,above]{no (start /\\game\_over / won)} (idle);
  \draw[flow] (c1) -- (mv); \draw[flow] (mv) -- (cw); \draw[flow] (cw) -- (cp);
  \draw[flow] (cp) -- (pw); \draw[flow] (pw) -- (c2); \draw[flow] (c2) -- (fx);
  \draw[flow] (fx) -- (win);
  \draw[flow] (win) -- node[lbl,above]{yes} (won);
  \draw[flow] (win) -- node[lbl,right]{no} (over);
  \draw[flow] (over) -- node[lbl,right]{yes} (last);
  \draw[flow] (last) -- node[lbl,right]{yes} (go);
  \draw[flow] (last) -- node[lbl,above]{no} (miss);
  \draw[flow] (over.west) -- node[lbl,above]{no} ++(-2.2,0) |- (idle.south);
  \draw[flow] (miss.east) -- ++(0.8,0) |- (top.east);
  \draw[flow] (idle.north) -- ++(0,0.6) -- ++(0,1) -| (top.west);
\end{tikzpicture}}
\caption{One frame. Steps 1 to 7 run only while playing. The red-bordered step~6 is a
duplicate of step~1 (Section~\ref{sec:double}); the green path is the win condition.
All paths return to the top of the loop.}
\label{fig:loop}
\end{figure}

\begin{honest}{The loop has no frame-rate limiter}
There is no \code{time.sleep()} and no clock anywhere in the loop. It runs as fast as the
CPU allows, so the game's actual speed depends entirely on the machine it runs on: the ball
crosses the screen faster on a fast computer. Real engines fix this by either sleeping to a
target frame rate or scaling movement by elapsed time (``delta time''). Neither is done here.
The constants \code{BASE\_SPEED = 1.05} and \code{MAX\_SPEED = 2.6} are therefore tuned to
one particular machine and mean nothing in absolute terms.

Worth noting: the README's ``Setup'' section claims the game imports \code{time}. No module
in the project imports \code{time}. That import does not exist, which is consistent with
there being no frame limiter.
\end{honest}

\subsection{The state machine}

Physics must not run on the start screen or after a game ends, but the window must stay
responsive so it can still receive the \code{space} or \code{r} keypress. A single string
variable, \code{game\_state}, handles this.

\begin{figure}[H]
\centering
\begin{tikzpicture}[node distance=22mm]
  \node[box] (start) {\textbf{start}\\\scriptsize title screen, no physics};
  \node[box, right=of start] (playing) {\textbf{playing}\\\scriptsize full frame loop};
  \node[box, above right=12mm and 20mm of playing] (won) {\textbf{won}\\\scriptsize ``YOU WIN!''};
  \node[box, below right=12mm and 20mm of playing] (over) {\textbf{game\_over}\\\scriptsize ``GAME OVER''};

  \draw[flow] (start) -- node[lbl,above]{\code{space}\\\scriptsize\code{start\_game()}} (playing);
  \draw[flow] (playing) -- node[lbl,above left=-1mm and -3mm]{all blocks\\destroyed} (won);
  \draw[flow] (playing) -- node[lbl,below left=-1mm and -3mm]{last life\\lost} (over);
  \draw[flow] (won.north) to[out=100,in=60, looseness=1.4] node[lbl,above]{\code{r} \ \scriptsize\code{restart\_game()}} (start.north);
  \draw[flow] (over.south) to[out=-100,in=-60, looseness=1.2] node[lbl,below]{\code{r} \ \scriptsize\code{restart\_game()}} (start.south);
\end{tikzpicture}
\caption{Four states. Note the guards: \code{start\_game()} only acts if the state is
\code{start}, and \code{restart\_game()} only if it is \code{game\_over} or \code{won}.}
\end{figure}

\begin{keyidea}{Why the guards matter}
Keyboard callbacks in turtle fire whenever the key is pressed, regardless of what the game
is doing. Without the \code{if game\_state == "start"} guard on line 67, mashing
\code{space} mid-rally would do nothing harmful here, but mashing \code{r} would silently
rebuild the block wall underneath a live ball. The guards make the two callbacks no-ops in
every state where they do not apply, which is the correct way to handle input that arrives
asynchronously.
\end{keyidea}

\section{Collision handling}

There are three genuinely different collision problems in this game, and the project's
file split exists precisely because they are different shapes of problem. This section
takes each in turn.

\begin{figure}[H]
\centering
\resizebox{\textwidth}{!}{
\begin{tikzpicture}
  % play field
  \draw[draw=linegrey, fill=softgrey] (0,0) rectangle (10,7.5);
  \node[font=\scriptsize\itshape, text=black!60, anchor=south] at (5,7.6)
    {play field: 800 x 600 window, coordinates -400..400 by -300..300};

  % walls
  \draw[warnred, very thick] (0.35,0.35) rectangle (9.65,7.15);
  \node[lbl, text=warnred] at (5,6.85) {\code{x > 390} / \code{x < -390} / \code{y > 290} / \code{y < -290}};

  % blocks
  \foreach \i in {0,...,6}{
    \draw[fill=accent!30, draw=accent] (0.75+\i*1.2,5.85) rectangle (1.65+\i*1.2,6.2);
    \draw[fill=okgreen!25, draw=okgreen] (0.75+\i*1.2,5.35) rectangle (1.65+\i*1.2,5.7);
  }

  % ball and radius
  \fill[white, draw=inkblue] (4.2,3.6) circle (0.16);
  \draw[accent, dashed] (4.2,3.6) circle (0.62);
  \node[lbl] at (4.2,2.7) {radius 35};
  \draw[flow] (4.2,3.6) -- (5.0,4.4);
  \node[lbl] at (5.4,4.2) {\code{(dx, dy)}};

  % paddle
  \draw[fill=white, draw=inkblue, very thick] (3.6,0.95) rectangle (6.4,1.2);

  % game over line
  \draw[warnred, dashed, thick] (0.35,0.62) -- (9.65,0.62);
  \node[lbl, text=warnred, anchor=west] at (0.55,0.45) {\code{y < -280}: life lost};

  % annotations, anchored well clear of the field
  \node[font=\scriptsize, text=inkblue, align=left, anchor=west] at (10.9,5.8)
    {\textbf{(3) ball vs.\ block}\\\code{block.distance(ball) < 35}\\radius check, in \file{blocks.py}};
  \node[font=\scriptsize, text=warnred, align=left, anchor=west] at (10.9,3.4)
    {\textbf{(1) ball vs.\ wall}\\clamp then flip\\the matching axis,\\in \file{ball.py}};
  \node[font=\scriptsize, text=inkblue, align=left, anchor=west] at (10.9,1.1)
    {\textbf{(2) ball vs.\ paddle}\\\code{y < -240} \textbf{and}\\\code{px-80 < x < px+80}\\directional, in \file{ball.py}};

  \draw[dflow] (10.8,5.8) -- (9.1,5.8);
  \draw[dflow] (10.8,3.4) -- (9.75,3.4);
  \draw[dflow] (10.8,1.1) -- (6.5,1.1);
\end{tikzpicture}}
\caption{The three collision rules on one field, drawn to scale in the game's own
coordinate system. Note how close the bottom wall line (solid red, \code{y = -290}) sits
to the life-lost line (dashed red, \code{y = -280}). That gap is the subject of
Section~\ref{sec:deadwall}.}
\label{fig:collide}
\end{figure}

\subsection{Collision 1: ball versus wall}

\file{ball.py} lines 42 to 58. The window is 800 by 600, and turtle puts the origin at the
centre, so valid coordinates run from $-400$ to $400$ horizontally and $-300$ to $300$
vertically. The walls are checked at $\pm 390$ and $\pm 290$, a 10-unit inset that accounts
for the ball's own radius.

\begin{lstlisting}[language=Python]
def check_collision_wall(self):
    if self.xcor() > 390:
        self.reset_x(sign='+')   # setx(390): clamp back onto the boundary
        self.bounce_x()          # dx *= -1: reverse horizontal direction
    ...
\end{lstlisting}

\begin{keyidea}{Clamp, then flip}
Each rule does two things, and both are necessary. \code{bounce\_x()} reverses direction.
But because the ball moves in discrete jumps, the frame that detects the overshoot has
already placed the ball \emph{past} the wall. Flipping alone would send it back from
outside the field. \code{reset\_x()} snaps it exactly onto the boundary first, so the
bounce starts from a legal position. Without the clamp, a ball could get stuck vibrating
outside the wall: flip, still outside, flip again, forever.
\end{keyidea}

\subsection{Collision 2: ball versus paddle}

\file{ball.py} lines 60 to 66. This one is deliberately \emph{directional}:

\begin{lstlisting}[language=Python]
def check_collision_paddle(self, paddle):
    if (self.ycor() < -240) and (paddle.xcor() + 80 > self.xcor() > paddle.xcor() - 80):
        self.sety(-240)          # clamp, same reasoning as the wall
        self.dy *= -1            # bounce upward
        self.paddle_hits += 1
        if self.paddle_hits % HITS_PER_STEP == 0:
            self.speed_up()
\end{lstlisting}

Two conditions must both hold: the ball is below $y = -240$, and its $x$ is within 80 units
of the paddle's centre (the paddle is drawn with \code{stretch\_len=8}, and turtle's base
square is 20 units, so it is 160 wide, hence $\pm 80$).

\begin{honest}{The bounce is unconditional, not physical}
\code{self.dy *= -1} always sends the ball back the way it came vertically, and \code{dx}
is never touched. So the paddle is a perfect mirror: \emph{where} the ball strikes it has
no effect at all. In the original Breakout, hitting the paddle's edge angles the ball away,
which is the game's core skill expression. Here the player has no way to steer, which
removes the strategy from the game entirely. The only variation in the ball's angle across
a whole run comes from \code{reset\_after\_miss()} randomizing the horizontal sign.

This is a design gap, not a bug, and it is the single highest-value improvement available
to the project: deriving \code{dx} from the offset between the ball's $x$ and the paddle's
centre would be a handful of lines and would make it an actual game of skill.
\end{honest}

\begin{honest}{The paddle is a one-way mirror, and that is on purpose}
Because the check only fires when \code{ycor() < -240}, a ball travelling \emph{upward}
through the paddle's own y-band passes straight through it. That sounds like a bug and is
actually the fix for one. Without the direction condition, the frame after a successful
bounce would still find the ball inside the paddle's band, flip \code{dy} a second time,
and drive the ball back down into the paddle: it would stick and jitter instead of bouncing
away. The clamp to \code{sety(-240)} plus the strict \code{<} is what breaks that cycle.
The README describes this correctly.
\end{honest}

\subsection{Collision 3: ball versus block}

\file{blocks.py} lines 118 to 129. Turtle offers no rectangle intersection test, so the
project uses distance instead.

\begin{lstlisting}[language=Python]
def check_collision_ball(self, ball, scoreboard):
    for block in self.all_blocks:
        if block.distance(ball) < 35:
            scoreboard.increase_score()
            self.all_blocks.remove(block)
            x, y, color = block.xcor(), block.ycor(), block.break_color
            block.color("white")
            self.fading.append([block, FADE_FRAMES])
            self.spawn_burst(x, y, color)
            ball.dy *= -1
            return True
    return False
\end{lstlisting}

\begin{jargon}{Hitbox}
An invisible rectangle attached to a game object, used for collision maths instead of its
actual drawn pixels (which would be far more expensive to test). ``No hitboxes'' here means
the project substitutes a circle of fixed radius for every block's true rectangle.
\end{jargon}

\begin{keyidea}{Why \code{return True} on the first hit matters}
The loop returns the instant it finds a hit, rather than continuing through the rest of the
list. This is not just an optimisation. \code{self.all\_blocks.remove(block)} mutates the
list that the \code{for} loop is currently iterating over, which in Python causes the
iterator to silently \emph{skip} the next element. Returning immediately means the loop
never continues after the mutation, so the bug never gets a chance to manifest. It is
correct, but it is correct by accident: if someone later removed the \code{return} to allow
multiple blocks per frame, they would introduce a skipping bug that is genuinely hard to
spot.
\end{keyidea}

\begin{honest}{One radius for blocks of three different widths}
Blocks are 60, 80 or 100 units wide (\code{stretch\_len} of 3, 4 or 5, times turtle's
20-unit base square) and 20 tall. The collision test treats every one of them as a circle
of radius 35.

That is wrong in both directions at once. A 100-wide block's far corners sit about 51 units
from its centre, so the ball can visibly overlap a wide block's end and \emph{not} register
a hit. Meanwhile 35 is far larger than the block's 10-unit half-height, so the ball can
register a hit while still visibly above or below a block. The README frames this as a
tolerable approximation, which understates it: the error is roughly $\pm 45\%$ of a block's
half-width depending on which block. A bounding-box test is about four comparisons and
would be strictly cheaper than \code{distance()}, which computes a square root.
\end{honest}

\begin{honest}{Only \code{dy} flips, so the ball passes through vertical faces}
\code{ball.dy *= -1} assumes every block is struck from above or below. A ball arriving
almost horizontally along a row will bounce vertically anyway, which looks visibly wrong:
it continues sideways through the row, destroying blocks in a line while flipping up and
down. A correct implementation compares the overlap on each axis and flips whichever is
smaller. This is the standard fix and it is not implemented.
\end{honest}

\label{sec:double}
\begin{honest}{The double collision call is a workaround, and it does not work}
\file{main.py} calls \code{block.check\_collision\_ball()} at line 97 and again at line
112, in the same iteration, with nothing in between that could plausibly change the answer
except \code{ball.move()}. The README says it appears to be there to stop a fast ball
tunnelling past a block.

It cannot achieve that. Both calls test the ball's \emph{position at that instant}; neither
tests the \emph{path} it swept. The ball moves at most 2.6 units per frame against a
35-unit radius, so tunnelling is arithmetically impossible in the first place, and the
second call is pure cost. Worse, because each call can flip \code{dy} and destroy a block,
a ball sitting near two blocks can lose two blocks in a single frame and have \code{dy}
flipped \emph{twice}, which cancels out: it scores two points and carries on in its original
direction as if nothing had been hit. Removing the second call would be a strict improvement.
\end{honest}

\section{\file{ball.py}: motion and difficulty}

\subsection{Velocity}

\begin{jargon}{Velocity, \code{dx} and \code{dy}}
Speed plus direction. It is stored as two numbers: \code{dx} is how far the ball moves
horizontally per frame, \code{dy} vertically. The ``d'' is for ``delta'', meaning change.
Moving is then just addition, which is all \code{move()} does:
\code{goto(xcor() + dx, ycor() + dy)}. A negative \code{dx} means leftward. Bouncing is
multiplying one of them by $-1$.
\end{jargon}

Both start at \code{BASE\_SPEED = 1.05}, so the ball always launches at exactly 45 degrees
up and to the right. Since \code{dx} is only ever negated and \code{dy} only negated or
scaled by the same factor as \code{dx}, \textbf{the ball moves at 45 degrees for the entire
game}, in one of four directions. Confirmed by inspection: no code path ever sets \code{dx}
and \code{dy} to different magnitudes.

\subsection{The difficulty ramp}

\begin{lstlisting}[language=Python]
BASE_SPEED = 1.05 ; SPEED_STEP = 1.08 ; MAX_SPEED = 2.6 ; HITS_PER_STEP = 4

def speed_up(self):
    new_dx = self.dx * SPEED_STEP
    new_dy = self.dy * SPEED_STEP
    if abs(new_dx) <= MAX_SPEED:
        self.dx = new_dx
    if abs(new_dy) <= MAX_SPEED:
        self.dy = new_dy
\end{lstlisting}

Every 4th paddle return multiplies both components by 1.08, an 8\% increase, checked
against a cap. Starting at 1.05, eleven increases are accepted, reaching
$1.05 \times 1.08^{11} = 2.4482$ at the 44th paddle return. The twelfth, due at the 48th
return, would give $2.6441$, which fails the \code{<= MAX\_SPEED} test and is rejected. The
speed therefore never changes again, and because nothing here is random, the sequence of
speeds is fixed and identical in every game. (Verified by running the ramp.)

\begin{keyidea}{Multiplying preserves the direction}
Scaling by a positive number leaves the sign alone, so a ball travelling left
($\code{dx} = -1.05$) speeds up to $-1.134$, still leftward. This is why the ramp multiplies
rather than adding: adding a constant would reverse a leftward ball's direction as it
crossed zero. Checking \code{abs(new\_dx)} against the cap is the same reasoning applied to
the limit.
\end{keyidea}

\begin{honest}{The cap makes the ramp silently stop, mid-flight}
Because each component is guarded by its own \code{if}, and both always have equal
magnitude, they always hit the cap on the same call, so they never desynchronise in
practice. But the guard \emph{rejects} the offending increase rather than clamping it, so
the speed freezes at $2.4482$ and \code{MAX\_SPEED = 2.6} is never actually reached: the
ball tops out about 6\% below the limit the constant advertises. The constant names a speed
the ball can never travel at. Clamping with \code{min(new\_dx, MAX\_SPEED)} would honour it;
the current code does not. Confirmed by running the ramp to 80 paddle hits: the final speed
is 2.4482, last changed at hit 44.
\end{honest}

\subsection{The two resets}

\begin{center}
\begin{tabular}{@{}lll@{}}
\toprule
& \textbf{\code{reset\_after\_miss()}} & \textbf{\code{full\_reset()}} \\
\midrule
Called when      & a life is lost, lives remain & \code{r} pressed after game over/win \\
Position         & back to $(0,0)$              & back to $(0,0)$ \\
Speed            & \textbf{kept} (difficulty survives) & back to \code{BASE\_SPEED} \\
\code{paddle\_hits} & \textbf{kept}             & back to 0 \\
Direction        & horizontal sign randomized   & up and to the right \\
\bottomrule
\end{tabular}
\end{center}

The distinction is a real design decision and a good one: losing a life is a setback, not a
reprieve, so the game does not hand the player back an easy slow ball as a reward for
missing.

\begin{honest}{\code{reset\_after\_miss} sends the ball downward, into the pit}
\code{self.dy = speed\_y}, where \code{speed\_y = abs(self.dy)}. Taking the absolute value
forces \code{dy} positive, meaning \emph{upward}. That is deliberate and correct.

But only \code{dx}'s sign is randomized (\code{random.random() < 0.5}); \code{dy} is always
positive. The README says the ball is sent off in ``a new direction'', which is true of only
one of the two axes. In practice this is the right call, since a downward-launching ball
would fall straight past the paddle for a free second life loss, but the randomization is
half of what the comment in the code (``a fresh (randomized) direction'') implies.
\end{honest}

\section{\file{blocks.py}: building the wall}

\subsection{The packing algorithm}

This is the most intricate code in the project. \code{make\_all\_blocks()} loops calling
\code{make\_block()}, which picks a random width, then hands the block to
\code{place\_block()} to decide where it goes. Blocks are packed \textbf{right to left},
starting at $x = 400$, wrapping down a row when the next one will not fit, and stopping
when there is no vertical room left.

\begin{figure}[H]
\centering
\resizebox{0.95\textwidth}{!}{
\begin{tikzpicture}
  \draw[draw=linegrey, dashed] (0,0.6) rectangle (13,6.2);
  \node[lbl, anchor=south] at (0.5,6.25) {\code{x=-400}};
  \node[lbl, anchor=south] at (12.5,6.25) {\code{x=400}};

  % direction of packing
  \draw[dflow] (12.7,5.95) -- (1.1,5.95);
  \node[lbl, fill=white] at (6.9,5.95) {packed right to left};

  % row 1
  \foreach \x/\w in {11.2/1.8, 9.2/2.0, 7.4/1.8, 5.0/2.4, 3.2/1.8, 1.2/2.0}
    \draw[fill=warnred!25, draw=warnred] (\x,5.3) rectangle (\x+\w-0.06,5.7);
  \node[lbl, anchor=east] at (-0.1,5.5) {\code{y=240}};

  % row 2
  \foreach \x/\w in {11.4/1.6, 9.0/2.4, 7.2/1.8, 5.2/2.0, 2.8/2.4, 0.9/1.9}
    \draw[fill=orange!30, draw=orange!80!black] (\x,4.0) rectangle (\x+\w-0.06,4.4);
  \node[lbl, anchor=east] at (-0.1,4.2) {\code{y=220}};

  % wrap: leftmost of row 1 runs out of room, so the next block starts at the
  % right of the row below. Routed through the empty band under row 2.
  \draw[flow] (1.2,5.5) -- (0.4,5.5) -- (0.4,3.7) -- (12.7,3.7) -- (12.7,3.94);
  \node[lbl, fill=white] at (6.5,3.7) {next block would cross \code{x=-400}: wrap to \code{y -= 20}};

  % row 3 partial
  \foreach \x/\w in {11.2/1.8, 9.4/1.8, 7.0/2.4}
    \draw[fill=yellow!55!orange!45, draw=olive] (\x,2.7) rectangle (\x+\w-0.06,3.1);
  \node[lbl, anchor=east] at (-0.1,2.9) {\code{y=200}};
  \node[lbl] at (3.6,2.9) {\ldots\ five rows total, down to \code{y=160}};

  \draw[warnred, thick, dashed] (0,1.6) -- (13,1.6);
  \node[lbl, text=warnred, fill=white, anchor=north] at (6.5,1.5)
    {next row would be \code{y=140}, which fails \code{y > LOWER\_YCOR}: stop};
\end{tikzpicture}}
\caption{Packing. Widths are random per block (60, 80 or 100 units), so rows do not line
up. Each row takes its colour from \code{ROW\_COLORS} by its $y$ position, which is what
makes the wall read as bands rather than a random speckle.}
\end{figure}

The wrap test, \file{blocks.py} lines 57 to 61, computes where the next block's centre would
land and compares it against a limit that already accounts for the block's own half-width:

\begin{lstlisting}[language=Python]
x = ((self.x_y_position[-1][0] - ((self.random_stretch_len[-2] * 20) // 2))
     - ((self.random_stretch_len[-1] * 20) // 2))
lim = (UPPER_XCOR - ((self.random_stretch_len[-1] * 20) // 2)) * -1
if x < lim:
    ...  # wrap to a new row
\end{lstlisting}

Read plainly: take the previous block's centre, step left by the previous block's
half-width (to reach its left edge), then step left again by the \emph{new} block's
half-width. That lands the new block flush against the old one. If that centre is further
left than \code{lim}, it will not fit, so wrap.

\begin{keyidea}{Why the code needs both \code{[-1]} and \code{[-2]}}
\code{random\_stretch\_len[-1]} is the block being placed right now; \code{[-2]} is the one
placed before it. The two half-widths are different numbers and both are needed, because
the gap between two block centres is the sum of their half-widths, not twice either one.
This is why the class keeps parallel lists at all rather than just tracking the last
position.
\end{keyidea}

\subsection{Block-break feedback}

Rather than a block simply vanishing, breaking one does three things: the block turns white
and is queued in \code{self.fading} for 3 frames, five particle turtles are spawned flying
outward, and only then does everything disappear.

\begin{lstlisting}[language=Python]
angle = random.uniform(0, 360)
dx = math.cos(math.radians(angle)) * 5
dy = math.sin(math.radians(angle)) * 5
self.particles.append([particle, dx, dy, BURST_LIFETIME])
\end{lstlisting}

\begin{jargon}{Sine, cosine and why they are here}
To scatter particles evenly in every direction you need to turn an angle into an
$(x, y)$ direction. Cosine of an angle gives the horizontal part, sine the vertical, each
between $-1$ and $1$; multiplying both by 5 sets the speed. Picking the angle at random
from 0 to 360 therefore picks a random direction at a fixed speed, which is exactly what a
burst is. \code{math.radians} converts degrees into the unit Python's trigonometry expects.
\end{jargon}

Note the iteration idiom in \code{update\_effects()}: \code{for entry in self.fading[:]}.
The \code{[:]} makes a shallow copy, so the loop iterates the copy while removing from the
original. This is the correct and explicit fix for the mutation-during-iteration hazard
that \code{check\_collision\_ball()} only dodges by luck. The same author wrote both; one is
right for the right reason.

\begin{keyidea}{Effects outlive the round on purpose}
\file{main.py} line 132 calls \code{block.update\_effects()} in the \code{else} branch too,
so a burst that was in flight when the last life was lost keeps animating on the game-over
screen instead of freezing mid-air. This is a small, deliberate touch and it is easy to
miss when reading the loop.
\end{keyidea}

\section{\file{paddle.py} and \file{scoreboard.py}}

\subsection{The paddle}

The simplest module, 34 lines. It binds Left and Right arrows to methods that step $\pm 20$
units, and clamps itself to $\pm 320$ each frame so it cannot leave the screen.

\begin{honest}{\code{onkeypress} without \code{onkeyrelease} means movement is not smooth}
The paddle moves 20 units per key \emph{event}, not per frame. Held-key repeat is handled by
the operating system, so the paddle's speed and its initial delay before repeating are the
OS keyboard settings, not the game's. The frame-based way to do this is to record the key as
held on press, clear it on release, and move by a fixed amount each frame inside the loop.
That would make paddle motion independent of keyboard settings, and it is a two-method change.
\end{honest}

\subsection{The scoreboard and persistence}

\begin{jargon}{Persistence}
Keeping data after the program exits. Variables live in memory and die with the process; to
survive, a value must be written to a file (or a database). Here the single persisted value
is the high score, in \file{highestScore.txt}.
\end{jargon}

The read happens at \textbf{module level}, at lines 3 to 8, outside any function or class:

\begin{lstlisting}[language=Python]
try:
    score = int(open('highestScore.txt', 'r').read())
except FileNotFoundError:
    score = open('highestScore.txt', 'w').write(str(0))
except ValueError:
    score = 0
\end{lstlisting}

\begin{honest}{Verified bug: a first run starts the high score at 1, not 0}
\code{file.write()} returns \textbf{the number of characters it wrote}, not the value
written. Writing the string \code{"0"} writes one character, so \code{write()} returns
\code{1}, and the \code{FileNotFoundError} branch assigns \code{score = 1}.

On any machine with no \file{highestScore.txt}, the game therefore opens claiming a high
score of 1. The player must destroy two blocks before their score is recognised as beating
it. The README states this branch means the game ``just starts the high score at 0''. It
does not.

Confirmed by running the module with the file absent:
\begin{lstlisting}
H1 module-level `score` on first run (no file) = 1
\end{lstlisting}
The fix is one line: assign \code{score = 0} and write the file as a separate statement.
The bug is invisible after the first run, which is exactly why it survived.
\end{honest}

\begin{honest}{Import has a side effect, and it is a write}
Reading the file at module level means \code{import scoreboard} touches the disk, and in
the missing-file case \emph{creates} a file. An import should not do that: it makes the
module impossible to test without a filesystem, impossible to reuse with a different save
location, and it means the read happens once per process no matter how many
\code{Scoreboard} objects exist. The README's ``What I'd do differently'' identifies this
correctly and it is the right call.

There is a second, quieter consequence. Because the file is read once at import, the value
in \code{Scoreboard.highScore} is a snapshot from process start. Two copies of the game
running at once would each overwrite the other's high score with a stale value on exit.
\end{honest}

\begin{honest}{\code{increase\_score} increments the high score instead of assigning it}
\begin{lstlisting}[language=Python]
self.score += 1
if self.score > self.highScore:
    self.highScore += 1
\end{lstlisting}
It happens to be correct, because \code{score} only ever rises by 1, so the moment it
overtakes \code{highScore} the two are equal and stay locked together. But the intent is
``the high score is the best score'', and the code says ``add one to the high score''.
\code{self.highScore = self.score} expresses the intent and would survive any future change
that awarded more than one point per block. As written, awarding 5 points per block would
silently break the high score.
\end{honest}

\begin{honest}{\code{Scoreboard.reset()} is dead code}
Lines 56 to 58 define \code{reset()}, which calls \code{save\_high\_score()} then
\code{new\_round()}. Nothing calls it: \file{main.py} calls \code{save\_high\_score()} and
\code{new\_round()} separately, at the two moments it actually wants each. Confirmed by
grep across the project. It is a leftover from an earlier structure and should be deleted.
\end{honest}

\section{Honest findings}
\label{sec:honest}

The boxes above cover each issue where it arises. This section collects the ones found by
\textbf{running} the code rather than reading it, since those are the ones an interviewer
could surface that reading alone would not.

\subsection{The ghost block}

\begin{honest}{Verified bug: an invisible block sits at the origin and scores a free point}
Trace \code{make\_block()} carefully:

\begin{lstlisting}[language=Python]
def make_block(self):
    block = Turtle()                    # a new turtle is born at (0, 0)
    ...
    self.place_block(block)             # may return early WITHOUT calling goto()
    row = int(round((UPPER_YCOR - block.ycor()) / 20))
    ...
    self.all_blocks.append(block)       # appended REGARDLESS
\end{lstlisting}

When the wall runs out of vertical room, \code{place\_block()} sets
\code{can\_make\_more = False} and returns \emph{before} calling \code{block.goto()}. The
block is left at the turtle default position, $(0,0)$, the dead centre of the play field.
\code{make\_block()} then appends it to \code{all\_blocks} anyway. \code{make\_all\_blocks()}
calls \code{hideturtle()} on it, which makes it invisible but does \textbf{not} remove it
from the list.

The consequences are all real:
\begin{itemize}
  \item \code{check\_collision\_ball()} iterates \code{all\_blocks}, so the ghost is a live
        collider. The ball starts at $(0,0)$: distance 0. On the \textbf{first frame of
        every game}, the player scores a free point off an invisible block, \code{dy} flips
        so the ball launches \emph{downward} rather than upward, and a particle burst fires
        at the centre of the screen from nowhere.
  \item \code{len(block.all\_blocks) == 0} is the win condition, so the win requires
        destroying the ghost too. It is destroyed on frame one, so the win still works, but
        by luck.
  \item The row calculation for the ghost is \code{round((240 - 0) / 20) = 12}, and
        \code{ROW\_COLORS[12 \% 7]} is \code{'steel blue'}. The burst at the centre is
        steel blue, a colour no real row ever uses.
\end{itemize}

Confirmed by running the real \code{Block} and \code{Ball} classes:
\begin{lstlisting}
H2 len(all_blocks) after make_all_blocks = 50
H2 last block pos = (0.00,0.00) isvisible = False color = steel blue
H4 ball start pos = (0.00,0.00)  score before = 0
H4 collision on first frame = True  score after = 1  ball.dy = -1.05
\end{lstlisting}

Every visible symptom of this bug (the centre burst, the ball starting downward, the score
reading 1) looks like a feature if you have never read the code. The fix is one line: do not
append the block when \code{can\_make\_more} went false, or \code{pop()} it in
\code{make\_all\_blocks()} instead of merely hiding it.
\end{honest}

\subsection{The dead bottom wall}
\label{sec:deadwall}

\begin{honest}{Verified dead code: the bottom wall bounce can never execute}
\file{ball.py} bounces off the bottom wall at \code{ycor() < -290}. \code{check\_game\_over()}
fires at \code{ycor() < -280}. In the loop (Figure~\ref{fig:loop}), the ball moves, then the
wall is checked, then, at the end of the same iteration, game-over is checked. The ball's
maximum speed is 2.6 units per frame, so it cannot cross the 10-unit gap between $-280$ and
$-290$ within one frame. Game-over always fires first, and the ball is teleported away.

Driving the ball straight down at maximum speed through main's exact per-frame order:
\begin{lstlisting}
game_over fired at y = -280.6 on frame 30
lowest y ever reached before game_over = -280.6
bottom-wall bounce threshold is y < -290 ; reachable? False
\end{lstlisting}

The bottom wall is unreachable code. This is harmless (it is a floor under a floor) but it
is misleading: the README claims ``the ball bounces off all four walls'', and it bounces off
three. The code reads as though there is a safety net that does not exist.
\end{honest}

\subsection{Two unreachable colours}

\begin{honest}{\code{ROW\_COLORS} has 7 entries and the wall can only ever have 5 rows}
Rows are placed at $y = 240, 220, 200, 180, 160$; the next would be $140$, which fails the
\code{y > LOWER\_YCOR} test since \code{LOWER\_YCOR = 140}. So there are always exactly five
rows, using \code{ROW\_COLORS[0..4]}. \code{'steel blue'} appears only on the invisible ghost
block, and \code{'medium purple'} never appears at all.

The comment on line 10, ``cycled if there are more rows than colors'', describes a case the
geometry makes impossible. The modulo in \code{ROW\_COLORS[row \% len(ROW\_COLORS)]} is
defensive code guarding against nothing. Confirmed:
\begin{lstlisting}
H3 distinct y rows present = [0, 160, 180, 200, 220, 240]
H3 colors in use = ['gold', 'medium sea green', 'orange', 'steel blue',
                    'tomato', 'yellow green']
\end{lstlisting}
(The \code{0} row and \code{'steel blue'} in that output are the ghost block.)
\end{honest}

\subsection{Documentation drift}

\begin{honest}{Three README claims the code contradicts}
\begin{enumerate}
  \item \textbf{``a 279-line turtle game''} (``What I'd do differently''). Measured:
        \code{wc -l} gives 475 lines across the five modules (134 + 147 + 102 + 58 + 34).
        The figure is stale by roughly 40\%.
  \item \textbf{``uses only \ldots turtle \ldots and time''} (``Setup''). Nothing imports
        \code{time}. The actual imports are \code{turtle}, \code{random} and \code{math}.
  \item \textbf{``a corrupted file doesn't crash the game, it just starts the high score at
        0''} (``Challenges''). It starts at 0 for a corrupted file (the \code{ValueError}
        branch) but at \textbf{1} for a missing file, as proven above.
\end{enumerate}
These are cheap to fix and expensive to be caught on.
\end{honest}

\subsection{Verification note}

\begin{keyidea}{What was actually run, and what was not}
Every claim in this section was checked by importing the project's real modules against a
real display and driving them, working on copies in a scratch directory so the repository's
own \file{highestScore.txt} was never read or written. The probes replicate \file{main.py}'s
exact per-frame call order.

The full interactive game was \textbf{not} played end to end: \file{main.py} never returns
until the window closes, and neither the paddle's key bindings nor the visual appearance of
the particle bursts were exercised. Claims about how things \emph{look} (the flash, the
burst, the colour bands) are read from the code and are marked as such rather than seen.
\end{keyidea}

\section{What the project gets right}

The honest findings are numerous, so it is worth being equally precise about the parts that
are genuinely well done, because they are the parts worth defending:

\begin{itemize}
  \item \textbf{The file split earns itself.} Each collision rule is a different shape of
        problem and each lives in the module that owns the relevant object. This is the
        right decomposition and it is the reason the bugs above are findable at all.
  \item \textbf{\code{tracer(0)} plus a manual \code{update()}} is the correct and
        non-obvious way to get frame-based rendering out of turtle.
  \item \textbf{The two-reset design} (keep difficulty after a miss, drop it on restart) is
        a real game-feel decision, not an accident.
  \item \textbf{Clamp-then-flip} in the wall and paddle collisions shows the author hit the
        stuck-vibrating bug and fixed it properly.
  \item \textbf{The \code{[:]} copy} in \code{update\_effects()} shows awareness of the
        mutation-during-iteration hazard.
  \item \textbf{Catching \code{Terminator}} so closing the window exits cleanly rather than
        dumping a traceback is a detail most projects at this size skip.
  \item \textbf{The README's ``What I'd do differently'' is largely correct}, and identifies
        the module-level file read, the distance check and the double collision call
        unprompted. The drift documented above is in the details, not the judgement.
\end{itemize}

\begin{keyidea}{The honest summary}
This is a competent exercise that does the thing it set out to do: it builds a working game
loop, a state machine and three kinds of collision detection without an engine, and the
author can explain why each file exists. It also ships a free point from an invisible block
on frame one, a high score that starts at 1, a wall of three walls, and a paddle that cannot
be aimed with. Every one of those is a small fix and a good story about how it was found.
\end{keyidea}

\end{document}
